\chapter{2000年新课标卷（理）}
\section{单选题}
\begin{problem}
设集合 \(A\) 和 \(B\) 都是坐标平面上的点集 \(\{(x, y) \mid x \in \R,\ y \in \R\}\)，映射 \(f: A \to B\) 把集合 \(A\) 中的元素 \((x, y)\) 映射成集合 \(B\) 中的元素 \((x + y, x - y)\)，则在映射 \(f\) 下，象 \((2, 1)\) 的原象是（\quad）

\choices
    {\((3, 1)\)}
    {\(\left(\frac{3}{2}, \frac{1}{2}\right)\)}
    {\(\left(\frac{3}{2}, -\frac{1}{2}\right)\)}
    {\((1, 3)\)}
\end{problem}

\begin{answer}
B
\end{answer}

\begin{problem}
在复平面内，把复数 \(3 - \sqrt{3}\i\) 对应的向量按顺时针方向旋转 \(\frac{\pi}{3}\)，所得向量对应的复数是（\quad）

\choices
    {\(2\sqrt{3}\)}
    {\(-2\sqrt{3}\i\)}
    {\(\sqrt{3} - 3\i\)}
    {\(3 + \sqrt{3}\i\)}
\end{problem}

\begin{answer}
B
\end{answer}

\begin{problem}
一个长方体共一顶点的三个面的面积分别是 \(\sqrt{2}\)，\(\sqrt{3}\)，\(\sqrt{6}\)，这个长方体对角线的长是（\quad）

\choices
    {\(2\sqrt{3}\)}
    {\(3\sqrt{2}\)}
    {\(6\)}
    {\(\sqrt{6}\)}
\end{problem}

\begin{answer}
D
\end{answer}

\begin{problem}
设 \(\bs{a}\)，\(\bs{b}\)，\(\bs{c}\) 是任意的非零平面向量，且相互不共线，则有下列命题：

\begin{circlelist}
\item \((\bs{a} \cdot \bs{b})\bs{c} - (\bs{c} \cdot \bs{a})\bs{b} = 0\)；
\item \(|\bs{a}| - |\bs{b}| < |\bs{a} - \bs{b}|\)；
\item \((\bs{b} \cdot \bs{c})\bs{a} - (\bs{c} \cdot \bs{a})\bs{b}\) 不与 \(\bs{c}\) 垂直；
\item \((3\bs{a} + 2\bs{b}) \cdot (3\bs{a} - 2\bs{b}) = 9|\bs{a}|^{2} - 4|\bs{b}|^{2}\).
\end{circlelist}

其中是真命题的有（\quad）

\choices
    {\circled{1}\circled{2}}
    {\circled{2}\circled{3}}
    {\circled{3}\circled{4}}
    {\circled{2}\circled{4}}
\end{problem}

\begin{answer}
D
\end{answer}

\begin{problem}
函数 \( y = -x\cos x \) 的部分图象是（\quad）

\choices
    {\choicebitmap[width=0.15\paperwidth]{img/2000/new_curriculum_science/q5_optA.png}}
    {\choicebitmap[width=0.15\paperwidth]{img/2000/new_curriculum_science/q5_optB.png}}
    {\choicebitmap[width=0.15\paperwidth]{img/2000/new_curriculum_science/q5_optC.png}}
    {\choicebitmap[width=0.15\paperwidth]{img/2000/new_curriculum_science/q5_optD.png}}
\end{problem}

\begin{answer}
D
\end{answer}

\begin{problem}
《中华人民共和国个人所得税法》规定，公民全月工资、薪金所得不超过 \(800\text{ 元}\) 的部分不必纳税，超过 \(800\text{ 元}\) 的部分为全月应纳税所得额，此项税款按下表分段累进计算：

\begin{center}
\begin{tabular}{|c|c|}
\hline
全月应纳税所得额 & 税率 \\
\hline
不超过 \(500\text{ 元}\) 的部分 & \(5\%\) \\
\hline
超过 \(500\text{ 元}\) 至 \(2000\text{ 元}\) 的部分 & \(10\%\) \\
\hline
超过 \(2000\text{ 元}\) 至 \(5000\text{ 元}\) 的部分 & \(15\%\) \\
\hline
\(\cdots\) & \(\cdots\) \\
\hline
\end{tabular}
\end{center}

某人一月份交纳此项税款 \(26.78\text{ 元}\)，则他的当月工资、薪金所得介于（\quad）

\choices
    {\(800\text{ 元} - 900\text{ 元}\)}
    {\(900\text{ 元} - 1200\text{ 元}\)}
    {\(1200\text{ 元} - 1500\text{ 元}\)}
    {\(1500\text{ 元} - 2800\text{ 元}\)}
\end{problem}

\begin{answer}
C
\end{answer}

\begin{problem}
若 \(a > b > 1\)，\(P = \sqrt{\lg a \cdot \lg b}\)，\(Q = \frac{1}{2} (\lg a + \lg b)\)，\(R = \lg \left(\frac{a + b}{2}\right)\)，则（\quad）

\choices
    {\(R < P < Q\)}
    {\(P < Q < R\)}
    {\(Q < P < R\)}
    {\(P < R < Q\)}
\end{problem}

\begin{answer}
B
\end{answer}

\begin{problem}
图中阴影部分的面积是（\quad）

\bitmapfigure[max width=0.65\linewidth]{img/2000/new_curriculum_science/q8_fig1.png}

\choices
    {\(2\sqrt{3}\)}
    {\(9 - 2\sqrt{3}\)}
    {\(\frac{32}{3}\)}
    {\(\frac{35}{3}\)}
\end{problem}

\begin{answer}
C
\end{answer}

\begin{problem}
一个圆柱的侧面展开图是一个正方形，这个圆柱的全面积与侧面积的比是（\quad）

\choices
    {\(\frac{1 + 2\pi}{2\pi}\)}
    {\(\frac{1 + 4\pi}{4\pi}\)}
    {\(\frac{1 + 2\pi}{\pi}\)}
    {\(\frac{1 + 4\pi}{2\pi}\)}
\end{problem}

\begin{answer}
A
\end{answer}

\begin{problem}
过原点的直线与圆 \(x^{2} + y^{2} + 4x + 3 = 0\) 相切，若切点在第三象限，则该直线的方程是（\quad）

\choices
    {\(y = \sqrt{3}x\)}
    {\(y = -\sqrt{3}x\)}
    {\(y = \frac{\sqrt{3}}{3}x\)}
    {\(y = -\frac{\sqrt{3}}{3}x\)}
\end{problem}

\begin{answer}
C
\end{answer}

\begin{problem}
过抛物线 \(y = ax^{2}\)（\(a > 0\)）的焦点 \(F\) 作一直线交抛物线于 \(P\)，\(Q\) 两点，若线段 \(PF\) 与 \(FQ\) 的长分别是 \(p\)，\(q\)，则 \(\frac{1}{p} + \frac{1}{q}\) 等于（\quad）

\choices
    {\(2a\)}
    {\(\frac{1}{2a}\)}
    {\(4a\)}
    {\(\frac{4}{a}\)}
\end{problem}

\begin{answer}
C
\end{answer}

\begin{problem}
如图，\(OA\) 是圆锥底面中心 \(O\) 到母线的垂线，\(OA\) 绕轴旋转一周所得曲面将圆锥分成体积相等的两部分，则母线与轴的夹角为（\quad）

\bitmapinlinefigure[0.65\linewidth][]{img/2000/new_curriculum_science/q12_fig1.png}

\choices
    {\(\arccos\frac{1}{\sqrt[3]{2}}\)}
    {\(\arccos\frac{1}{2}\)}
    {\(\arccos\frac{1}{\sqrt{2}}\)}
    {\(\arccos\frac{1}{\sqrt[4]{2}}\)}
\end{problem}

\begin{answer}
D
\end{answer}

\section{填空题}
\begin{problem}
某厂生产电子元件，其产品的次品率为 \(5\%\)，现从一批产品中任意地连续取出 \(2\) 件，其中次品 \(\xi\) 的概率分布是：

\begin{center}
\begin{tabular}{|c|c|c|c|}
\hline
\(\xi\) & \(0\) & \(1\) & \(2\) \\
\hline
\(p\) &\fillinblank{}&\fillinblank{}&\fillinblank{}\\
\hline
\end{tabular}
\end{center}
\end{problem}

\begin{answer}
\(0.9025\)，\(0.095\)，\(0.0025\)
\end{answer}

\begin{problem}
椭圆 \(\frac{x^{2}}{9} + \frac{y^{2}}{4} = 1\) 的焦点 \(F_{1}\)，\(F_{2}\)，点 \(P\) 为其上的动点，当 \(\angle F_{1}PF_{2}\) 为钝角时，点 \(P\) 横坐标的取值范围是\fillinblank{}.
\end{problem}

\begin{answer}
\(\left(-\frac{3}{\sqrt{5}},\frac{3}{\sqrt{5}}\right)\)
\end{answer}

\begin{problem}
设 \(\{a_{n}\}\) 是首项为 \(1\) 的正项数列，且 \((n + 1)a_{n + 1}^{2} - na_{n}^{2} + a_{n + 1}a_{n} = 0\)（\(n = 1\)，\(2\)，\(3\)，\(\cdots\)），则它的通项公式是 \(a_{n}=\)\fillinblank{}.
\end{problem}

\begin{answer}
\(\frac{1}{n}\)
\end{answer}

\begin{problem}
如图，\(E\)，\(F\) 分别为正方体面 \(ADD_{1}A_{1}\)，面 \(BCC_{1}B_{1}\) 的中心，则四边形 \(BFD_{1}E\) 在该正方体的面上的射影可能是\fillinblank{}.

（要求：把可能的图序号都填上）

\begingroup
\leftskip=0pt\relax
\rightskip=0pt\relax
\examfiguregroup[float=inline]{img/2000/new_curriculum_science/q16_fig1.png}{%
    \vbox{%
        \hbox to \linewidth{%
            \hfil\bitmapinclude[max width=0.82\linewidth]{img/2000/new_curriculum_science/q16_fig1.png}\hfil
        }%
        \vskip 0.4\baselineskip
        \hbox to \linewidth{%
            \makebox[0.22\linewidth][c]{\bitmapinclude[max width=0.20\linewidth]{img/2000/new_curriculum_science/q16_optA.png}}\hfill
            \makebox[0.22\linewidth][c]{\bitmapinclude[max width=0.20\linewidth]{img/2000/new_curriculum_science/q16_optB.png}}\hfill
            \makebox[0.22\linewidth][c]{\bitmapinclude[max width=0.20\linewidth]{img/2000/new_curriculum_science/q16_optC.png}}\hfill
            \makebox[0.22\linewidth][c]{\bitmapinclude[max width=0.20\linewidth]{img/2000/new_curriculum_science/q16_optD.png}}%
        }%
    }%
}
\endgroup
\end{problem}

\begin{answer}
\(\circled{2}\circled{3}\)
\end{answer}

\section{解答题}
\begin{problem}
甲、乙二人参加普法知识竞答，共有 \(10\) 个不同的题目，其中选择题 \(6\) 个，判断题 \(4\) 个，甲、乙二人依次各抽一题.

\begin{enumerate}
\item 甲抽到选择题，乙抽到判断题的概率是多少？
\item 甲、乙二人中至少有一人抽到选择题的概率是多少？
\end{enumerate}
\end{problem}

\begin{solution}
\begin{enumerate}
\item 甲抽到选择题的概率为 \(\frac{6}{10}\). 甲抽走一个选择题后，还剩下 \(9\) 个题目，其中判断题仍有 \(4\) 个，所以乙抽到判断题的概率为 \(\frac{4}{9}\). 因此，甲抽到选择题且乙抽到判断题的概率为
\[
\frac{6}{10}\cdot\frac{4}{9}=\frac{4}{15}
\]
\item 先求甲、乙二人都抽到判断题的概率. 甲抽到判断题的概率为 \(\frac{4}{10}\)，甲抽走一个判断题后，乙抽到判断题的概率为 \(\frac{3}{9}\)，故二人都抽到判断题的概率为
\[
\frac{4}{10}\cdot\frac{3}{9}=\frac{2}{15}
\]
至少有一人抽到选择题是“不是二人都抽到判断题”，所以所求概率为
\[
1-\frac{2}{15}=\frac{13}{15}
\]
\end{enumerate}
\end{solution}

\section{选做题：解答题}

\begin{problem}[甲]
如图，直三棱柱 \(ABC - A_{1}B_{1}C_{1}\)，底面 \(\triangle ABC\) 中，\(CA = CB = 1\)，\(\angle BCA = 90^{\circ}\)，棱 \(AA_{1} = 2\)，\(M\)，\(N\) 分别是 \(A_{1}B_{1}\)，\(A_{1}A\) 的中点.

\begin{enumerate}
\item 求 \(\overline{BN}\) 的长；
\item 求 \(\cos\langle\overrightarrow{BA_{1}},\overrightarrow{CB_{1}}\rangle\) 的值；
\item 求证 \(A_{1}B \perp C_{1}M\).
\end{enumerate}

\bitmapinlinefigure[0.75\linewidth][]{img/2000/new_curriculum_science/q18_fig1.png}
\end{problem}

\begin{solution}
建立空间直角坐标系，取 \(C\) 为原点，分别以 \(CA\)，\(CB\)，\(CC_{1}\) 的方向为三个坐标轴的正方向.
由题意可得 \(A=(1,0,0)\)，\(B=(0,1,0)\)，\(C=(0,0,0)\)，\(A_{1}=(1,0,2)\)，\(B_{1}=(0,1,2)\)，\(C_{1}=(0,0,2)\).
又因为 \(M\)，\(N\) 分别是 \(A_{1}B_{1}\)，\(A_{1}A\) 的中点，所以
\(M=\left(\frac{1}{2},\frac{1}{2},2\right)\)，\(N=(1,0,1)\).

\begin{enumerate}
\item \(\overrightarrow{BN}=(1,-1,1)\)，所以
\[
|BN|=\sqrt{1^{2}+(-1)^{2}+1^{2}}=\sqrt{3}
\]
\item \(\overrightarrow{BA_{1}}=(1,-1,2)\)，\(\overrightarrow{CB_{1}}=(0,1,2)\)，从而
\[
\overrightarrow{BA_{1}}\cdot\overrightarrow{CB_{1}}=1\cdot0+(-1)\cdot1+2\cdot2=3
\]
且 \(|\overrightarrow{BA_{1}}|=\sqrt{1^{2}+(-1)^{2}+2^{2}}=\sqrt{6}\)，
\(|\overrightarrow{CB_{1}}|=\sqrt{0^{2}+1^{2}+2^{2}}=\sqrt{5}\).
因此
\[
\cos\langle\overrightarrow{BA_{1}},\overrightarrow{CB_{1}}\rangle=\frac{\overrightarrow{BA_{1}}\cdot\overrightarrow{CB_{1}}}{|\overrightarrow{BA_{1}}|\,|\overrightarrow{CB_{1}}|}=\frac{3}{\sqrt{30}}=\frac{\sqrt{30}}{10}
\]
\item \(\overrightarrow{A_{1}B}=(-1,1,-2)\)，\(\overrightarrow{C_{1}M}=\left(\frac{1}{2},\frac{1}{2},0\right)\)，所以
\[
\overrightarrow{A_{1}B}\cdot\overrightarrow{C_{1}M}=-\frac{1}{2}+\frac{1}{2}+0=0
\]
故 \(A_{1}B\perp C_{1}M\).
\end{enumerate}
\end{solution}

\reuseproblemnumber
\begin{problem}[乙]
如图，已知平行六面体 \(ABCD - A_{1}B_{1}C_{1}D_{1}\) 的底面 \(ABCD\) 为菱形，且 \(\angle C_{1}CB = \angle C_{1}CD = \angle BCD\).

\begin{enumerate}
\item 证明：\(C_{1}C \perp BD\)；
\item 当 \(\dfrac{CD}{CC_{1}}\) 的值为多少时，能使 \(A_{1}C \perp\) 平面 \(C_{1}BD\)？请给出证明.
\end{enumerate}

\bitmapinlinefigure[0.7\linewidth][]{img/2000/new_curriculum_science/q19_fig1.png}
\end{problem}

\begin{solution}
设 \(C\) 为原点，记 \(\bs{u}=\overrightarrow{CB}\)，\(\bs{v}=\overrightarrow{CD}\)，\(\bs{w}=\overrightarrow{CC_{1}}\)，并设 \(|\bs{u}|=|\bs{v}|=s\)，\(|\bs{w}|=h\)，\(\angle BCD=\theta\).
由题设 \(\angle C_{1}CB=\angle C_{1}CD=\angle BCD=\theta\)，可得 \(\bs{u}\cdot\bs{v}=s^{2}\cos\theta\)，\(\bs{w}\cdot\bs{u}=hs\cos\theta\)，\(\bs{w}\cdot\bs{v}=hs\cos\theta\).

\begin{enumerate}
\item 因为 \(\overrightarrow{BD}=\bs{v}-\bs{u}\)，所以
\[
\overrightarrow{CC_{1}}\cdot\overrightarrow{BD}=\bs{w}\cdot(\bs{v}-\bs{u})=hs\cos\theta-hs\cos\theta=0
\]
故 \(C_{1}C\perp BD\).
\item 由平行六面体的向量关系，\(\overrightarrow{CA}=\bs{u}+\bs{v}\)，\(\overrightarrow{CA_{1}}=\bs{u}+\bs{v}+\bs{w}\). 因而直线 \(A_{1}C\) 的方向向量可取 \(\bs{u}+\bs{v}+\bs{w}\). 先有
\[
(\bs{u}+\bs{v}+\bs{w})\cdot(\bs{v}-\bs{u})=0
\]
所以 \(A_{1}C\perp BD\). 又
\[
(\bs{u}+\bs{v}+\bs{w})\cdot(\bs{u}-\bs{w})=s^{2}+s^{2}\cos\theta-hs\cos\theta-h^{2}
\]
由对称性，\((\bs{u}+\bs{v}+\bs{w})\cdot(\bs{v}-\bs{w})\) 也等于上式. 因此，\(A_{1}C\perp\) 平面 \(C_{1}BD\) 的充要条件是
\[
s^{2}+s^{2}\cos\theta-hs\cos\theta-h^{2}=0
\]
令 \(r=\frac{h}{s}>0\)，上式除以 \(s^{2}\) 得
\[
r^{2}+r\cos\theta-(1+\cos\theta)=0
\]
即
\[
(r-1)(r+1+\cos\theta)=0
\]
由于 \(0<\theta<\pi\)，有 \(r+1+\cos\theta>0\)，所以 \(r=1\). 因此
\[
\frac{CD}{CC_{1}}=\frac{s}{h}=1
\]
当 \(r=1\) 时，上述三个内积均为零，故该条件也充分.
\end{enumerate}
\end{solution}

\section{解答题}

\begin{problem}
设 \(\{a_{n}\}\) 为等差数列，\(S_{n}\) 为数列 \(\{a_{n}\}\) 的前 \(n\) 项和，已知 \(S_{7} = 7\)，\(S_{15} = 75\)，\(T_{n}\) 为数列 \(\left\{\frac{S_{n}}{n}\right\}\) 的前 \(n\) 项和，求 \(T_{n}\).
\end{problem}

\begin{solution}
设等差数列 \(\{a_{n}\}\) 的首项为 \(a_{1}\)，公差为 \(d\). 由等差数列前 \(n\) 项和公式
\[
S_{n}=\frac{n}{2}\left(2a_{1}+(n-1)d\right)
\]
以及已知条件，得到
\[
\begin{cases}
2a_{1}+6d=2,\\
2a_{1}+14d=10
\end{cases}
\]
两式相减得 \(8d=8\)，所以 \(d=1\)，代入第一式得 \(a_{1}=-2\). 因此
\[
S_{n}=\frac{n}{2}\left(-4+n-1\right)=\frac{n(n-5)}{2}
\]
从而
\[
\frac{S_{k}}{k}=\frac{k-5}{2}
\]
所以
\[
T_{n}=\sum_{k=1}^{n}\frac{S_{k}}{k}=\sum_{k=1}^{n}\frac{k-5}{2}=\frac{1}{2}\left(\frac{n(n+1)}{2}-5n\right)=\frac{n(n-9)}{4}
\]
故 \(T_{n}=\frac{n(n-9)}{4}\).
\end{solution}

\begin{problem}
设函数 \(f(x) = \sqrt{x^{2} + 1} - ax\)，其中 \(a > 0\).

\begin{enumerate}
\item 解不等式 \(f(x) \leqslant 1\)；
\item 证明：当 \(a \geqslant 1\) 时，函数 \(f(x)\) 在区间 \([0, +\infty)\) 上是单调函数.
\end{enumerate}
\end{problem}

\begin{solution}
\begin{enumerate}
\item 因为 \(\sqrt{x^{2}+1}\geqslant0\)，所以不等式 \(f(x)\leqslant1\) 成立时必须有 \(ax+1\geqslant0\). 因此原不等式等价于 \(\sqrt{x^{2}+1}\leqslant ax+1\)，并且 \(x\geqslant-\frac{1}{a}\).
两边平方，得
\[
x^{2}+1\leqslant a^{2}x^{2}+2ax+1
\]
移项并因式分解，得
\[
x\left((1-a^{2})x-2a\right)\leqslant0
\]
下面分情况讨论.

当 \(0<a<1\) 时，\(1-a^{2}>0\)，上式的两个根为 \(0\) 和 \(\frac{2a}{1-a^{2}}\)，故
\[
0\leqslant x\leqslant\frac{2a}{1-a^{2}}
\]

当 \(a=1\) 时，上式化为 \(-2x\leqslant0\)，结合 \(x\geqslant-1\)，得
\[
x\geqslant0
\]

当 \(a>1\) 时，上式的两个根为 \(-\frac{2a}{a^{2}-1}\) 和 \(0\)，故由二次不等式得
\[
x\leqslant-\frac{2a}{a^{2}-1}\quad\text{或}\quad x\geqslant0
\]
又因为
\[
\frac{2a}{a^{2}-1}>\frac{1}{a}
\]
所以 \(-\frac{2a}{a^{2}-1}< -\frac{1}{a}\)，前一个区间与 \(x\geqslant-\frac{1}{a}\) 没有交集，仍得 \(x\geqslant0\). 综上，不等式的解集为
\[
\begin{cases}
\left[0,\frac{2a}{1-a^{2}}\right]&0<a<1\\
[0,+\infty)&a\geqslant1
\end{cases}
\]
\item 设 \(0\leqslant x_{1}<x_{2}\). 则
\[
\begin{aligned}
f(x_{2})-f(x_{1})
&=\sqrt{x_{2}^{2}+1}-\sqrt{x_{1}^{2}+1}-a(x_{2}-x_{1})\\
&=(x_{2}-x_{1})\left(\frac{x_{1}+x_{2}}{\sqrt{x_{2}^{2}+1}+\sqrt{x_{1}^{2}+1}}-a\right)
\end{aligned}
\]
由于 \(x_{1},x_{2}\geqslant0\)，有
\[
\sqrt{x_{2}^{2}+1}+\sqrt{x_{1}^{2}+1}>x_{1}+x_{2}
\]
所以括号内的第一项小于 \(1\)，而 \(a\geqslant1\)，从而 \(f(x_{2})-f(x_{1})<0\). 因此，函数 \(f(x)\) 在 \([0,+\infty)\) 上单调递减，当然是单调函数.
\end{enumerate}
\end{solution}

\begin{problem}
用总长 \(14.8\mathrm{~m}\) 的钢条制成一个长方体容器的框架，如果所制做容器的底面的一边比另一边长 \(0.5\mathrm{~m}\)，那么高为多少时容器的容积最大？并求出它的最大容积.
\end{problem}

\begin{solution}
设容器底面较短的一边长为 \(x\mathrm{~m}\)，则另一边长为 \((x+0.5)\mathrm{~m}\)，高为 \(h\mathrm{~m}\). 长方体框架共有四条长、四条宽和四条高，所以
\[
4\left(x+x+0.5+h\right)=14.8
\]
即
\[
h=3.2-2x
\]
由边长为正，得 \(0<x<1.6\). 容器的容积为
\[
V(x)=x(x+0.5)(3.2-2x)=-2x^{3}+2.2x^{2}+1.6x
\]
求导得
\[
V'(x)=-6x^{2}+4.4x+1.6=-6(x-1)\left(x+\frac{4}{15}\right)
\]
当 \(0<x<1\) 时，\(V'(x)>0\)；当 \(1<x<1.6\) 时，\(V'(x)<0\). 所以 \(V(x)\) 在 \(x=1\) 处取得最大值. 此时
\[
h=3.2-2=1.2
\]
且
\[
V_{\max}=1\times1.5\times1.2=1.8
\]
故当高为 \(1.2\mathrm{~m}\) 时容积最大，最大容积为 \(1.8\mathrm{~m}^{3}\).
\end{solution}

\begin{problem}
如图，已知梯形 \(ABCD\) 中 \(|AB| = 2|CD|\)，点 \(E\) 分有向线段 \(\overrightarrow{AC}\) 所成的比为 \(\lambda\)，双曲线过 \(C\)，\(D\)，\(E\) 三点，且以 \(A\)，\(B\) 为焦点，当 \(\frac{2}{3} \leqslant \lambda \leqslant \frac{3}{4}\) 时，求双曲线离心率 \(e\) 的取值范围.

\bitmapinlinefigure[0.55\linewidth][]{img/2000/new_curriculum_science/q23_fig1.png}
\end{problem}

\begin{solution}
取 \(CD=1\)，则 \(AB=2\). 建立直角坐标系，令 \(A=(0,0)\)，\(B=(2,0)\)，并设 \(D=(u,h)\)，\(C=(u+1,h)\)，其中 \(h>0\). 对于直线 \(y=h\) 上横坐标为 \(x\) 的点，记两个焦点距离之差为
\[
\Phi(x)=\sqrt{x^{2}+h^{2}}-\sqrt{(x-2)^{2}+h^{2}}
\]
由于 \(h>0\)，\(g'(x)=\frac{h^{2}}{(x^{2}+h^{2})^{3/2}}>0\)，所以函数 \(g(x)=\frac{x}{\sqrt{x^{2}+h^{2}}}\) 严格递增，且 \(\Phi'(x)=g(x)-g(x-2)>0\)，从而 \(\Phi(x)\) 严格递增，并且
\[
\Phi(2-x)=-\Phi(x)
\]
双曲线上各点到两焦点的距离之差的绝对值相等. 因为 \(C\)，\(D\) 的横坐标不同，\(\Phi(u)\) 与 \(\Phi(u+1)\) 不可能相等，所以
\[
\Phi(u)=-\Phi(u+1)=\Phi(1-u)
\]
由 \(\Phi\) 严格递增，得 \(u=1-u\)，即 \(u=\frac{1}{2}\). 因此 \(D=\left(\frac{1}{2},h\right)\)，\(C=\left(\frac{3}{2},h\right)\).
令 \(p=AD=BC=\sqrt{h^{2}+\frac{1}{4}}\)，\(r=AC=BD=\sqrt{h^{2}+\frac{9}{4}}\)，\(d=r-p\).
则 \(C\) 到两焦点的距离之差为 \(d\)，而双曲线的焦半距为 \(1\)，故
\[
2a=d
\]
设 \(t=\frac{AE}{AC}\). 由 \(\frac{AE}{EC}=\lambda\)，得
\[
t=\frac{\lambda}{1+\lambda}
\]
且 \(\frac{2}{3}\leqslant\lambda\leqslant\frac{3}{4}\) 时 \(0<t<\frac{1}{2}\). 因而 \(E=\left(\frac{3}{2}t,th\right)\) 在左支上，双曲线定义给出
\[
EB-AE=d
\]
又 \(AE=tr\)，所以
\[
EB=tr+d
\]
由坐标计算，
\[
EB^{2}=\left(2-\frac{3}{2}t\right)^{2}+t^{2}h^{2}
\]
注意到 \(r^{2}-p^{2}=2\)，且 \(d=r-p\)，所以
\[
r+p=\frac{2}{d}
\]
所以 \(r=\frac{1}{2}\left(d+\frac{2}{d}\right)\)，\(p=\frac{1}{2}\left(\frac{2}{d}-d\right)\).
从而 \(h^{2}=p^{2}-\frac{1}{4}\). 将 \(EB=tr+d\) 两边平方并代入上式，得
\[
\begin{gathered}
0=\left(2-\frac{3}{2}t\right)^{2}+t^{2}\left(p^{2}-\frac{1}{4}\right)-(tr+d)^{2}\\
{}=4-6t-2trd-d^{2}=4-8t-(1+t)d^{2}
\end{gathered}
\]
故
\[
d^{2}=\frac{4(1-2t)}{1+t}
\]
双曲线的离心率为
\[
e=\frac{c}{a}=\frac{1}{d/2}=\frac{2}{d}
\]
所以
\[
e^{2}=\frac{4}{d^{2}}=\frac{1+t}{1-2t}=\frac{1+2\lambda}{1-\lambda}
\]
函数 \(\frac{1+2\lambda}{1-\lambda}\) 在 \(\left[\frac{2}{3},\frac{3}{4}\right]\) 上递增，故
\[
7=\frac{1+2\cdot\frac{2}{3}}{1-\frac{2}{3}}\leqslant e^{2}\leqslant\frac{1+2\cdot\frac{3}{4}}{1-\frac{3}{4}}=10
\]
由于 \(e>0\)，双曲线离心率的取值范围为
\[
e\in[\sqrt{7},\sqrt{10}]
\]
\end{solution}
